\chapter{Functions of Several Variables and Partial Derivatives}
\label{ch32:several-variables-partial-derivatives}

Up through Chapter~31, where we approximated curves with Taylor polynomials, every function we
studied was one-dimensional: it took a single number $x$ and returned a single number $y$, and
its graph was a curve in the plane with a tangent line. Most quantities in science and
engineering, however, depend on several inputs at once: temperature in a room depends on three
spatial coordinates and time, aircraft costs depend on hundreds of parameters, and the loss
function of a neural network depends on millions of weights.

Here we step from the line onto the plane, and eventually into space. A function will now take
\emph{several} numbers as input and return one. Its graph will no
longer be a curve but a \emph{surface}. The single derivative will split into several
\emph{partial derivatives}, one for each input direction. And the simple intuition ``a
limit just means you can get as close as you like'' will turn out to hide a genuine new
subtlety, because in the plane there are infinitely many directions from which to approach a
point. Almost everything you already know transfers -- but it transfers with a twist, and
the twist is exactly where the new geometry lives.

\section{Functions of Two and Three Variables}
\label{ch32:sec-functions}

A function of two variables is a rule that assigns to each ordered pair $(x,y)$ a single
real number $z = f(x,y)$. The set of pairs we are allowed to feed in is the \emph{domain}
(now a region in the plane $\R^2$, not an interval on the line), and the set of outputs is
the \emph{range}. Nothing about the \emph{idea} of a function has changed: it is still a
machine, still single-valued, still ``one output per input.'' What changed is that the input
slot is wider. It now takes a point of the plane.

\begin{focusbox}
A real-valued function of two variables is a map $f : D \To \R$ where $D \subseteq \R^2$.
We write $z = f(x,y)$. For three variables, $f : D \To \R$ with $D \subseteq \R^3$, written
$w = f(x,y,z)$. The output is always a \emph{single number}; only the input has grown.
\end{focusbox}

Here are functions you should hold in your head for the rest of the chapter.

\begin{example}[A handful of surfaces]
\label{ch32:ex-zoo}
\begin{itemize}
  \item $f(x,y) = x^2 + y^2$. The height above $(x,y)$ is its squared distance from the
  origin. This is a bowl (a \emph{paraboloid}) opening upward.
  \item $f(x,y) = \sqrt{1 - x^2 - y^2}$, with domain the disk $x^2+y^2 \le 1$. This is the
  upper half of the unit sphere: at the origin the height is $1$, and it falls to $0$ at the
  rim.
  \item $f(x,y) = x^2 - y^2$. A \emph{saddle}: it curves up as you walk along the $x$-axis
  and down as you walk along the $y$-axis.
  \item $T(x,y,z) = \dfrac{1}{x^2+y^2+z^2}$, a function of three variables -- think of it as
  the temperature near a hot point source at the origin, blazing at the center and cooling
  with distance.
\end{itemize}
\end{example}

Notice the domain is now a genuinely two-dimensional object. For $\sqrt{1-x^2-y^2}$ we must
have $x^2+y^2 \le 1$, so the domain is a filled disk. For $\ln(x+y)$ we need $x+y>0$, the
open half-plane above the line $y=-x$. Sketching the domain is the first thing to do with any
new multivariable function, because it tells you \emph{where the surface lives}.

\begin{intuitionbox}
Think of $z = f(x,y)$ as the elevation of rolling terrain. The pair $(x,y)$ is your position on
a map -- east and north -- and $f(x,y)$ is the \emph{altitude} at that spot. The graph is the
mountain topography itself: hills, valleys, ridges, passes. Every word we are about to invent --
level curve, partial derivative, tangent plane -- has a direct hiking interpretation. When in
doubt, go for a walk on the mountain.
\end{intuitionbox}

\subsection*{Three variables and why you cannot draw them}

A function $w=f(x,y,z)$ of three variables would have a graph living in four-dimensional
space (three for input, one for output). You cannot draw that, and neither can anyone else.
This is not a defect of your imagination; it is a fact about paper. The remedy is to study
three-variable functions through their \emph{level sets} (below) and through slices, holding
some inputs fixed. We will lean on two-variable functions for pictures and let the
three-variable case ride along on the same algebra, because every definition we give extends
to $n$ variables by simply writing more commas.

\section{Graphs, Level Curves, and Surfaces}
\label{ch32:sec-graphs}

There are two complementary ways to see a function of two variables. The first is the
\emph{graph}: lift each domain point $(x,y)$ to the height $z=f(x,y)$ and look at the surface
that results. The second is the \emph{contour map}: stay flat on the domain and draw the
curves where the height is constant, exactly as a topographic map encodes a mountain on a flat
sheet. Cartographers solved the ``draw a 3D surface on 2D paper'' problem centuries before
calculus needed it, and their solution -- contour lines -- is one of the most useful pictures
in all of multivariable calculus.

\begin{definition}[Level curve, level surface]
\label{ch32:def-level}
For a function $f(x,y)$ and a constant $c$, the \emph{level curve} (or contour) at level $c$
is the set of domain points where the function equals $c$:
\[
   \set{ (x,y) : f(x,y) = c }.
\]
For $f(x,y,z)$, the analogous set $\set{(x,y,z): f(x,y,z)=c}$ is a \emph{level surface}.
\end{definition}

\begin{figure}[H]
\centering
\begin{tikzpicture}
\begin{axis}[
title style={font=\small},
  width=8.4cm, height=7.2cm,
  view={40}{30},
  xlabel={$x$}, ylabel={$y$}, zlabel={$z$},
  domain=-2:2, y domain=-2:2,
  samples=28, samples y=28,
  colormap/viridis,
  title={The paraboloid $z=x^2+y^2$},
]
\addplot3[surf, shader=interp, opacity=0.92] {x^2 + y^2};
\end{axis}
\end{tikzpicture}
\hspace{0.3cm}
\begin{tikzpicture}
\begin{axis}[
title style={font=\small},
  width=6.6cm, height=7.2cm,
  xlabel={$x$}, ylabel={$y$},
  axis equal image,
  title={Its level curves $x^2+y^2=c$},
  domain=-2:2, y domain=-2:2,
]
\foreach \c in {1,2,3,4} {
  \edef\temp{\noexpand\addplot[curveblue, thick, samples=80, domain=0:360]
    ({sqrt(\c)*cos(x)}, {sqrt(\c)*sin(x)});}
  \temp
}
\end{axis}
\end{tikzpicture}
\caption{Left: the bowl $z=x^2+y^2$ as a surface. Right: its contour map. Each circle is a
level curve $x^2+y^2=c$; the circles are evenly spaced in height but get \emph{closer together}
as you move outward, because the bowl gets steeper. Crowded contours mean steep terrain.}
\label{ch32:fig-paraboloid}
\end{figure}

Read Figure~\ref{ch32:fig-paraboloid} the way a hiker reads a trail map. Each circle is a
fixed altitude. Where the circles bunch up, the ground is steep; where they spread apart, it
is gentle. For the bowl $x^2+y^2=c$, the level curve at height $c$ is a circle of radius
$\sqrt{c}$. The radii $\sqrt{1},\sqrt{2},\sqrt{3},\dots$ grow more and more slowly, so equal
steps in height correspond to ever-smaller steps outward -- which is precisely the contours
crowding together. The contour spacing \emph{is} the steepness, encoded in flatland.

\begin{happeningbox}
A level curve is a \emph{horizontal slice} of the graph, dropped straight down onto the
domain. Slice the surface $z=f(x,y)$ with the plane $z=c$; the intersection is a curve sitting
at height $c$ in space; its shadow on the floor is the level curve. So the contour map is the
collection of all these shadows, one per height. Nothing is lost: from a dense enough family
of contours you can reconstruct the whole surface. This is why a topographic map, though flat,
tells you everything about the mountain.
\end{happeningbox}

\begin{example}[Reading a saddle from its contours]
\label{ch32:ex-saddle}
Take $f(x,y)=x^2-y^2$. The level curve at $c$ is $x^2-y^2=c$. For $c>0$ these are hyperbolas
opening left-right; for $c<0$ they open up-down; and for $c=0$ we get $x^2=y^2$, i.e. the two
lines $y=\pm x$ crossing at the origin. That crossing pair is the signature of a saddle: the
contours \emph{intersect} themselves only at a saddle point, where the surface is going up in
one direction and down in another at the very same spot.
\end{example}

\begin{figure}[H]
\centering
\begin{tikzpicture}
\begin{axis}[
title style={font=\small},
  width=8.2cm, height=6.8cm,
  view={35}{28},
  xlabel={$x$}, ylabel={$y$}, zlabel={$z$},
  domain=-1.6:1.6, y domain=-1.6:1.6,
  samples=28, samples y=28,
  colormap/cool,
  title={The saddle $z=x^2-y^2$},
]
\addplot3[surf, shader=interp, opacity=0.9] {x^2 - y^2};
\end{axis}
\end{tikzpicture}
\hspace{0.2cm}
\begin{tikzpicture}
\begin{axis}[
title style={font=\small},
  width=6.6cm, height=6.8cm,
  xlabel={$x$}, ylabel={$y$},
  axis equal image,
  title={Contours of $x^2-y^2$},
  domain=-1.6:1.6, y domain=-1.6:1.6,
  samples=60, unbounded coords=jump,
]
\addplot[curveblue, thick, domain=-1.6:1.6] {x};
\addplot[curveblue, thick, domain=-1.6:1.6] {-x};
\foreach \c in {0.5,1.2} {
  \edef\temp{\noexpand\addplot[tangentred, domain=-1.6:1.6, samples=80]
    {sqrt(x*x - \c)};}
  \temp
  \edef\temp{\noexpand\addplot[tangentred, domain=-1.6:1.6, samples=80]
    {-sqrt(x*x - \c)};}
  \temp
}
\end{axis}
\end{tikzpicture}
\caption{The saddle and its contours. The blue lines $y=\pm x$ are the level set $z=0$ and
cross at the saddle point; the red hyperbolas are levels $c=0.5,1.2$. Self-crossing contours
are how a saddle announces itself on a flat map.}
\label{ch32:fig-saddle}
\end{figure}

\begin{walkbox}
On your next walk, find a mountain pass -- a saddle -- in your mind: the low point of a ridge
between two peaks. Stand on it. Forward and backward (toward the peaks) the ground rises;
left and right (down into the valleys) it falls. You are at a point that is simultaneously a
\emph{minimum} along one path and a \emph{maximum} along another. Hold that feeling; it is the
geometric heart of why finding maxima and minima in several variables is harder than in one.
\end{walkbox}

\begin{applicationbox}
Level sets are everywhere once you look. Weather maps draw \emph{isobars} (level curves of
pressure) and \emph{isotherms} (level curves of temperature). Economists draw
\emph{indifference curves}: level curves of a utility function $U(x,y)$, the combinations of two
goods that leave a consumer equally happy. In machine learning, the \emph{decision boundary}
of a classifier is a level set $f(\vec x)=0$ of the score function, separating one predicted
class from another. Whenever someone says ``the set of points that are equally good/hot/likely,''
they are naming a level set.
\end{applicationbox}

\section{Limits and Continuity in Several Variables}
\label{ch32:sec-limits}

Now the twist. We need a limit: what it means for $f(x,y)$ to approach a value $L$ as $(x,y)$
approaches a point $(a,b)$. The algebraic definition looks almost identical to the one-variable
version, but the geometry in the plane is much subtler.

On the line, ``$x \to a$'' has exactly two avenues of approach: from the left and from the
right. We made peace with that using one-sided limits, and a two-sided limit existed precisely
when both sides agreed. In the plane, ``$(x,y)\to(a,b)$'' allows approach from \emph{every
direction at once} -- straight in along the $x$-axis, in along the $y$-axis, in along any line,
in along a spiral, along a parabola. For the limit to exist, \emph{all} of these infinitely
many approaches must agree on the same value $L$. That is a far stronger demand.

\begin{definition}[Limit of a function of two variables]
\label{ch32:def-limit}
We say $\displaystyle\lim_{(x,y)\to(a,b)} f(x,y) = L$ if for every $\eps>0$ there is a
$\delta>0$ such that
\[
   0 < \sqrt{(x-a)^2 + (y-b)^2} < \delta
   \quad\Longrightarrow\quad
   \abs{f(x,y) - L} < \eps .
\]
In words: once $(x,y)$ is within distance $\delta$ of $(a,b)$ -- inside a small disk around
$(a,b)$, in any direction -- the output is within $\eps$ of $L$.
\end{definition}

Compare this letter for letter with the one-variable $\eps$--$\delta$ definition. The only
change is that $\abs{x-a}<\delta$, the length of an interval, has become
$\sqrt{(x-a)^2+(y-b)^2}<\delta$, the radius of a \emph{disk}. The interval became a disk
because the domain became two-dimensional. That single substitution carries all the new
subtlety with it, as shown in Figure~\ref{ch32:fig-limit-paths}.

\begin{figure}[H]
\centering
\begin{tikzpicture}
% Left panel: 1D approach along the line
\begin{scope}[shift={(0,0)}]
  \draw[thick, ->] (-0.3, 1.5) -- (5.8, 1.5) node[right, font=\small] {$x$};
  \fill[sourceblue!15] (1.3, 1.35) rectangle (4.3, 1.65);
  \draw[sourceblue!80, thick] (1.3, 1.35) -- (1.3, 1.65);
  \draw[sourceblue!80, thick] (4.3, 1.35) -- (4.3, 1.65);
  \fill[black] (2.8, 1.5) circle (2pt) node[below=6pt, font=\small] {$a$};
  \node[sourceblue!80!black, font=\scriptsize, above] at (1.3, 1.7) {$a-\delta$};
  \node[sourceblue!80!black, font=\scriptsize, above] at (4.3, 1.7) {$a+\delta$};
  \draw[sourceblue, very thick, -{Stealth[length=5pt]}] (1.5, 2.3) -- (2.6, 2.3)
    node[midway, above, font=\scriptsize] {from left};
  \draw[sourceblue, dashed] (1.5, 2.3) -- (1.5, 1.65);
  \draw[sourceblue, very thick, -{Stealth[length=5pt]}] (4.1, 2.3) -- (3.0, 2.3)
    node[midway, above, font=\scriptsize] {from right};
  \draw[sourceblue, dashed] (4.1, 2.3) -- (4.1, 1.65);
  \node[font=\footnotesize\bfseries, align=center] at (2.8, 0.4) {1D: Exactly two directions};
\end{scope}

% Right panel: 2D approach in the plane
\begin{scope}[shift={(7.4,0)}]
  \draw[gray!50, ->] (0, 1.5) -- (5.6, 1.5) node[right, font=\scriptsize, black] {$x$};
  \draw[gray!50, ->] (2.8, 0.2) -- (2.8, 3.4) node[above, font=\scriptsize, black] {$y$};
  \fill[sourceblue!12] (2.8, 1.8) circle (1.3);
  \draw[sourceblue!75, thick, dashed] (2.8, 1.8) circle (1.3);
  \fill[black] (2.8, 1.8) circle (2pt) node[below left=3pt, font=\scriptsize] {$(a,b)$};
  \draw[gray!70, -{Stealth[length=3.5pt]}] (2.8, 1.8) -- (3.72, 2.72)
    node[midway, above left, font=\tiny, black] {$\delta$};
  \draw[sourceblue, thick, -{Stealth[length=4pt]}] (1.6, 1.8) -- (2.6, 1.8);
  \draw[sourceblue, thick, -{Stealth[length=4pt]}] (2.8, 0.6) -- (2.8, 1.6);
  \draw[targetgreen, thick, -{Stealth[length=4pt]}] (3.7, 0.9) -- (3.0, 1.6);
  \draw[quotientorange, thick, -{Stealth[length=4pt]}] (3.9, 2.2) to[out=160, in=30] (2.95, 1.95);
  \draw[structurepurple, thick, -{Stealth[length=4pt]}] (1.8, 2.6) to[out=-20, in=120] (2.65, 1.95);
  \node[font=\tiny, quotientorange, anchor=west, fill=white, inner sep=1pt] at (4.05, 2.25) {parabola};
  \node[font=\tiny, targetgreen, below right] at (3.5, 0.8) {diagonal};
  \node[font=\tiny, structurepurple, above left] at (2.0, 2.6) {curved path};
  \node[font=\footnotesize\bfseries, align=center] at (2.8, 0.0) {2D: Infinitely many paths};
\end{scope}
\end{tikzpicture}
\caption{The leap from one variable to several. Left: On the real line, approaching $a$ allows
only two avenues -- from the left and from the right -- so a two-sided limit exists as soon as
those two agree. Right: In the plane, approaching $(a,b)$ requires entering a circular disk of
radius $\delta$ from every conceivable direction: along horizontal or vertical axes, along
diagonals, or along curved paths like parabolas. For the multivariable limit to exist, \emph{all}
of these infinitely many approaches must arrive at the exact same destination $L$.}
\label{ch32:fig-limit-paths}
\end{figure}

\begin{focusbox}
The quantity $\sqrt{(x-a)^2+(y-b)^2}$ is just the distance from $(x,y)$ to $(a,b)$. The
definition says: small distance forces small error, \emph{regardless of direction}. The
$\delta$-disk has no preferred direction, which is exactly why the limit must come out the
same no matter how you approach.
\end{focusbox}

\subsection*{The two-path test: how multivariable limits fail}

Here is the practical consequence, and it is the single most important computational fact in
this section. If you can find \emph{two} paths into $(a,b)$ that give two \emph{different}
limiting values, then the limit does not exist. One disagreement among the infinitely many
directions is enough to destroy the limit. Agreement along finitely many paths
\emph{never} proves a limit exists --- it only fails to disprove it.

\begin{example}[The classic non-limit]
\label{ch32:ex-twopath}
Consider
\[
   f(x,y) = \frac{xy}{x^2+y^2}, \qquad (x,y)\neq(0,0),
\]
and ask whether the limit exists as $(x,y)\to(0,0)$.

\textbf{Approach along the $x$-axis} ($y=0$): then $f(x,0) = \dfrac{x\cdot 0}{x^2+0} = 0$ for
all $x\neq 0$. So along this path $f \to 0$.

\textbf{Approach along the $y$-axis} ($x=0$): then $f(0,y) = 0$ likewise. Still $0$. So far so
innocent -- if these were the only two paths, we might be fooled.

\textbf{Approach along the diagonal} $y=x$: then
\[
   f(x,x) = \frac{x\cdot x}{x^2 + x^2} = \frac{x^2}{2x^2} = \frac12
   \quad\text{for all } x\neq 0 .
\]
Along the diagonal the function is \emph{constantly} $\tfrac12$, no matter how close we get to
the origin. The axes say $0$; the diagonal says $\tfrac12$. They disagree. Therefore
\[
   \lim_{(x,y)\to(0,0)} \frac{xy}{x^2+y^2} \quad\text{does not exist.}
\]
\end{example}

\begin{figure}[H]
\centering
\begin{tikzpicture}
\begin{axis}[
title style={font=\small},
  width=8.6cm, height=6.6cm,
  view={45}{25},
  xlabel={$x$}, ylabel={$y$}, zlabel={$z$},
  domain=-1:1, y domain=-1:1,
  samples=40, samples y=40,
  colormap/viridis,
  title={$z=\dfrac{xy}{x^2+y^2}$ near the origin},
  zmin=-0.6, zmax=0.6,
]
\addplot3[surf, shader=interp, opacity=0.95] {x*y/(x^2 + y^2 + 0.0008)};
\end{axis}
\end{tikzpicture}
\caption{The surface $z=xy/(x^2+y^2)$. Along the axes it sits at height $0$; along the diagonals
it rides at $\pm\tfrac12$. The surface is pleated right at the origin -- there is no single
height the function settles toward, so the limit fails.}
\label{ch32:fig-twopath}
\end{figure}

There is an even cleaner way to see the failure, straight from the contour map. Along the ray
at angle $\theta$ from the origin the function $xy/(x^2+y^2)$ is \emph{constant}: in polar form
it equals $\tfrac12\sin 2\theta$, which depends on $\theta$ but not on $r$. So every ray into
the origin carries its own fixed value, and the values disagree from one ray to the next.
A function whose level sets are rays through the point you are approaching cannot have a limit
there -- you can read its non-existence right off Figure~\ref{ch32:fig-twopath-contour}.

\begin{figure}[H]
\centering
\begin{tikzpicture}
\begin{axis}[
title style={font=\small},
  width=8.4cm, height=7.0cm,
  xlabel={$x$}, ylabel={$y$},
  axis equal image,
  domain=-2:2, y domain=-2:2,
  title={Level rays of $xy/(x^2+y^2)$},
  xmin=-2.1, xmax=2.1, ymin=-2.1, ymax=2.1,
]
% rays through the origin, each a level set; value v = (1/2) sin 2theta
\foreach \ang/\val in {0/0, 30/0.43, 45/0.5, 60/0.43, 90/0, 120/-0.43, 135/-0.5, 150/-0.43} {
  \edef\temp{\noexpand\addplot[curveblue, thick, samples=2, domain=0:2]
    ({x*cos(\ang)}, {x*sin(\ang)});
    \noexpand\addplot[curveblue, thick, samples=2, domain=0:2]
    ({-x*cos(\ang)}, {-x*sin(\ang)});}
  \temp
}
\addplot[tangentred, thick, samples=2, domain=0:2] ({x*cos(45)},{x*sin(45)});
\addplot[tangentred, thick, samples=2, domain=0:2] ({-x*cos(45)},{-x*sin(45)});
\node[tangentred, font=\footnotesize, fill=white, inner sep=1pt] at (axis cs:1.45,1.65) {$f=1/2$};
\node[black, font=\footnotesize, fill=white, inner sep=1pt] at (axis cs:1.75,0.28) {$f=0$};
\addplot[mark=*, mark size=1.6pt, black] coordinates {(0,0)};
\end{axis}
\end{tikzpicture}
\caption{The level sets of $xy/(x^2+y^2)$ are straight rays through the origin, because the
value $\tfrac12\sin 2\theta$ depends only on the direction $\theta$. The diagonal $y=x$ (red)
carries the constant value $\tfrac12$; the axes (horizontal/vertical) carry $0$. Different rays,
different values -- so no single limit can exist at the origin.}
\label{ch32:fig-twopath-contour}
\end{figure}

\begin{trapbox}
\textbf{Checking a few lines does not prove a limit exists.} Worse: there are functions whose
limit is the same along \emph{every straight line} through the origin and yet still fails,
because some \emph{curved} path disagrees. The standard cautionary tale is
$g(x,y)=\dfrac{x^2 y}{x^4+y^2}$. Along every line $y=mx$ you get
$g \to 0$. But along the parabola $y=x^2$,
\[
   g(x,x^2) = \frac{x^2\cdot x^2}{x^4 + x^4} = \frac{x^4}{2x^4} = \frac12 .
\]
So the line-limits all agree on $0$, the parabola says $\tfrac12$, and the limit does not
exist. Moral: matching paths can only \emph{disprove} a limit. To \emph{prove} one exists you
need an estimate that controls \emph{all} directions at once, like the squeeze below.
\end{trapbox}

\begin{whybox}
Why is the two-path test legitimate? Because the limit in
Definition~\ref{ch32:def-limit} demands a single $L$ that works inside the \emph{whole} disk.
Any particular path of approach lives inside that disk, so if the full limit is $L$, the limit
\emph{restricted to any path} must also be $L$ -- the path inherits the bound. Contrapositive:
if two paths produce two different restricted limits, no single $L$ can serve both, so the full
limit cannot exist. The test works because a path is a special case of ``approaching,'' and the
general statement, if true, forces every special case.
\end{whybox}

\subsection*{When the limit \emph{does} exist: the squeeze}

To prove a limit exists you must bound the error in a way that does not care about direction.
Polar coordinates are the natural tool, because they package ``distance from the center'' into a
single variable $r$.

\begin{example}[A limit that exists]
\label{ch32:ex-exists}
Show that $\displaystyle\lim_{(x,y)\to(0,0)} \frac{x^2 y}{x^2+y^2} = 0$.

Write $x=r\cos\theta$, $y=r\sin\theta$, so $x^2+y^2=r^2$ and $r\to 0$ \emph{means} the same as
$(x,y)\to(0,0)$, regardless of the angle $\theta$. Then
\[
   \abs{\frac{x^2 y}{x^2+y^2}}
   = \abs{\frac{r^2\cos^2\theta \cdot r\sin\theta}{r^2}}
   = r\,\abs{\cos^2\theta\,\sin\theta}
   \le r .
\]
The angle dropped out of the bound: whatever direction $\theta$ you came in along, the size is
at most $r$. As $r\to 0$ this is squeezed to $0$. So the limit is $0$, and now we mean it for
\emph{all} approaches at once. (Concretely, given $\eps>0$, take $\delta=\eps$; then distance
$r<\delta$ forces the value within $\eps$ of $0$.)
\end{example}

The contrast between Examples~\ref{ch32:ex-twopath} and~\ref{ch32:ex-exists} is the whole
lesson. In the bad case the angle $\theta$ survived ($f=\tfrac12\sin 2\theta$ in polar form,
genuinely depending on direction); in the good case the angle was crushed under a factor of
$r$ that goes to zero no matter what.

\begin{definition}[Continuity]
\label{ch32:def-continuity}
A function $f$ is \emph{continuous} at $(a,b)$ if $(a,b)$ is in the domain, the limit there
exists, and
\[
   \lim_{(x,y)\to(a,b)} f(x,y) = f(a,b).
\]
\end{definition}

The good news, which we will use without ceremony from here on: every function built from
$x$, $y$, constants, and the familiar one-variable functions by adding, multiplying, dividing
(where the denominator is nonzero), and composing is continuous wherever it is defined.
Polynomials in $x$ and $y$ are continuous everywhere; $\sin(xy)$, $e^{x-y}$, and
$\sqrt{1+x^2+y^2}$ are continuous everywhere; rational functions are continuous off their zero
denominators. The pathologies above all came from a $0/0$ at a single point. Away from such
points, life is calm.

\begin{marginnotebox}
The leap to several variables happened early and informally. Leonhard Euler, in the
mid-1700s, computed with functions of several variables and ``partial differences'' fluently,
long before anyone had formulated a rigorous $\eps$--$\delta$ definition of a limit even in
\emph{one} variable. The direction-dependent subtlety we just encountered was simply not on
the eighteenth-century radar; mathematicians computed slopes along coordinate lines and
produced successful models for fluid flow and vibrating sound. The disquieting counterexamples
-- such as functions continuous along every straight line yet failing to be continuous at a
point -- were discovered in the late nineteenth century by analysts seeking rock-solid
foundations. Those counterexamples are not perversity for its own sake; they are the geometric
guardrails that prove our $\delta$-disk definition is the right one.
\end{marginnotebox}

\section{Partial Derivatives: Slopes in the Coordinate Directions}
\label{ch32:sec-partials}

Now we differentiate. In one variable the derivative measured one slope: how fast $y$ changes
as $x$ changes. On a surface there is no single slope, because the answer depends on which way
you walk. Walk east and the ground may rise; walk north and it may fall. So we measure slope
\emph{one direction at a time}, starting with the two coordinate directions. That is the
partial derivative, and the idea behind it is almost embarrassingly simple.

\begin{intuitionbox}
To find how the surface tilts in the $x$-direction at a point, \emph{freeze $y$}. Holding
$y$ constant turns $f(x,y)$ back into an ordinary one-variable function of $x$ alone -- a
single curve carved out of the surface by the vertical plane $y=\text{const}$. Differentiate
that curve the old way. The result is the partial derivative with respect to $x$. To get the
$y$-direction, freeze $x$ instead. The multivariable derivative is nothing but the
one-variable derivative, applied along one axis with the other coordinate held still.
\end{intuitionbox}

\begin{definition}[Partial derivatives]
\label{ch32:def-partial}
The \emph{partial derivative of $f$ with respect to $x$} at $(a,b)$ is
\[
   \pdv{f}{x}(a,b) \;=\; \lim_{h\to 0} \frac{f(a+h,\,b) - f(a,b)}{h},
\]
and with respect to $y$,
\[
   \pdv{f}{y}(a,b) \;=\; \lim_{h\to 0} \frac{f(a,\,b+h) - f(a,b)}{h},
\]
whenever the limits exist. We also write $f_x$ and $f_y$ for these.
\end{definition}

Stare at the first formula. Only the \emph{first} slot moves: $a+h$ in the $x$-place, the
$y$-place glued at $b$. It is exactly the one-variable difference quotient, run on the
function $x \mapsto f(x,b)$. The curly $\partial$ ($d$'s rounder cousin) is the only new
symbol, and it means precisely ``derivative while the other variables are held fixed.''

\begin{happeningbox}
The notation $\partial$ is doing real work. The plain $d$ in $\dv{y}{x}$ promised that $y$
depended on \emph{nothing but} $x$. Here $f$ depends on both $x$ and $y$, so we cannot ask for
``the'' rate of change of $f$ -- there isn't one. The rounded $\partial$ is a flag: it says
\emph{``rate of change of $f$ as $x$ varies, with every other variable held constant.''} The
symbol $\partial f/\partial x$ is a single indivisible object; the ``$\partial f$'' and
``$\partial x$'' on their own are not quantities you may cancel. The notation was popularized
by Legendre and made standard by Jacobi in the 1840s, exactly so that nobody would mistake a
partial rate for a total one.
\end{happeningbox}

\begin{figure}[H]
\centering
\begin{tikzpicture}
% Left panel: 3D bowl with slicing plane y = 0.5
\begin{scope}[shift={(0,0)}]
\begin{axis}[
  width=7.8cm, height=6.8cm,
  view={42}{22},
  xlabel={$x$}, ylabel={$y$}, zlabel={$z$},
  domain=-1.6:1.6, y domain=-1.6:1.6,
  samples=24, samples y=24,
  colormap/blackwhite,
  title style={font=\small},
  title={3D Surface \& Slicing Plane $y=0.5$},
]
\addplot3[surf, opacity=0.22, shader=flat, draw=gray!40] {x^2 + y^2};
% Slicing plane y=0.5 as a transparent sheet
\addplot3[
  draw=sourceblue!60, fill=sourceblue!12, opacity=0.45,
  samples=2, samples y=2, domain=-1.6:1.6, y domain=0.5:0.501,
] {x^2 + y^2};
% Sliced curve on surface
\addplot3[curveblue, very thick, samples=60, domain=-1.6:1.6]
    ({x}, {0.5}, {x^2 + 0.25});
\addplot3[mark=*, mark size=1.8pt, black] coordinates {(1,0.5,1.25)};
\node[above right, font=\scriptsize] at (axis cs:1,0.5,1.25) {$(1,0.5,1.25)$};
\end{axis}
\end{scope}

% Right panel: The 2D cross-section z = x^2 + 0.25
\begin{scope}[shift={(8.2,0)}]
\begin{axis}[
  width=7.2cm, height=6.8cm,
  axis lines=middle,
  xlabel={$x$}, ylabel={$z$},
  xmin=-1.6, xmax=1.6, ymin=0, ymax=3.2,
  samples=100, domain=-1.5:1.5,
  title style={font=\small},
  title={2D Slice Curve: Slope is $f_x$},
  clip=true,
]
\addplot[curveblue, very thick] {x^2 + 0.25};
\addplot[tangentred, thick, dashed, domain=0.2:1.6] {1.25 + 2*(x - 1)};
\addplot[only marks, mark=*, mark size=1.8pt, black] coordinates {(1,1.25)};
\node[above left, font=\scriptsize] at (axis cs:1,1.25) {$(1,1.25)$};
\node[tangentred, font=\scriptsize, anchor=west, fill=white, inner sep=1pt] at (axis cs:0.4,0.3) {Tangent slope $f_x = 2$};
\end{axis}
\end{scope}
\end{tikzpicture}
\caption{The geometric essence of a partial derivative: slicing a surface with a coordinate plane.
Left: The vertical plane $y=0.5$ slices through the paraboloid $z=x^2+y^2$, holding $y$ fixed and
carving out a curve in space. Right: Viewed straight on in the slice plane, that curve is the ordinary
single-variable parabola $z(x) = x^2 + 0.25$. The partial derivative $\pdv{f}{x}(1,0.5) = 2$ is simply
the slope of the tangent line to this cross-sectional curve.}
\label{ch32:fig-slices}
\end{figure}

\begin{labbox}
Compute $f_x$ and $f_y$ for $f(x,y)=x^2+y^2$ at $(1,0.5)$ by hand, from the definition, before
reading on.

For $f_x$: hold $y=0.5$, so $f(x,0.5)=x^2+0.25$, whose derivative is $2x$; at $x=1$ that is
$\mathbf{2}$. For $f_y$: hold $x=1$, so $f(1,y)=1+y^2$, derivative $2y$; at $y=0.5$ that is
$\mathbf{1}$. So the bowl tilts upward twice as steeply in the $x$-direction (slope $2$) as in
the $y$-direction (slope $1$) at that point -- precisely the slope of $2$ captured by the
tangent line on the $y=0.5$ cross-section in Figure~\ref{ch32:fig-slices}.
\end{labbox}

\subsection*{Computing partials in practice}

You almost never go back to the limit definition. The trick is the one your intuition already
gave you: \textbf{to differentiate with respect to $x$, treat $y$ as a constant} and apply all
your one-variable rules. Then swap roles for $y$.

\begin{example}[Treat the other variable as a number]
\label{ch32:ex-compute}
Let $f(x,y) = x^3 y^2 + \sin(xy) + 4y$.

\emph{With respect to $x$} (so $y$ is a constant):
\[
   f_x = 3x^2 y^2 + \cos(xy)\cdot y + 0 = 3x^2 y^2 + y\cos(xy).
\]
The $x^3$ differentiated to $3x^2$ while $y^2$ rode along as a constant multiplier; the chain
rule on $\sin(xy)$ produced $\cos(xy)$ times the $x$-derivative of the inside,
$\partial(xy)/\partial x = y$; and the lone $4y$, a constant as far as $x$ is concerned, died.

\emph{With respect to $y$} (so $x$ is a constant):
\[
   f_y = x^3\cdot 2y + \cos(xy)\cdot x + 4 = 2x^3 y + x\cos(xy) + 4.
\]
Now $x^3$ was the constant multiplier, $y^2 \mapsto 2y$, the inside derivative is
$\partial(xy)/\partial y = x$, and $4y \mapsto 4$. Notice how each variable's computation is a
completely ordinary single-variable problem.
\end{example}

\begin{example}[Three variables, same idea]
\label{ch32:ex-three}
Let $g(x,y,z) = x^2 e^{yz}$. Then, freezing the other two each time,
\[
   g_x = 2x\,e^{yz}, \qquad
   g_y = x^2 z\,e^{yz}, \qquad
   g_z = x^2 y\,e^{yz}.
\]
For $g_y$ we held $x$ and $z$ fixed, so $e^{yz}$ has $z$ as the constant multiplying $y$ in the
exponent, giving the factor $z$; symmetrically for $g_z$. Three variables, three partials, each
a one-liner.
\end{example}

\begin{trapbox}
Two slips to guard against. First, do not forget the chain-rule factor when the other variable
sits \emph{inside} a composite. In $\sin(xy)$, differentiating with respect to $x$ gives
$y\cos(xy)$, not $\cos(xy)$ -- the inside $xy$ has $x$-derivative $y$. Second, do not kill a
term too eagerly: in $f_x$, a term like $4y$ vanishes (it is constant in $x$), but a term like
$4xy$ does \emph{not} -- it becomes $4y$. ``Constant in $x$'' means \emph{no $x$ appears at
all}, not merely that a $y$ appears.
\end{trapbox}

\begin{applicationbox}
Partial derivatives are the language of sensitivity. In thermodynamics, pressure, volume, and
temperature are tangled together, and physicists write $\left(\partial P/\partial V\right)_T$
-- ``how pressure responds to volume \emph{at fixed temperature}'' -- with the held variable in
a subscript precisely because the answer changes if you hold something else fixed instead. In
economics, $\partial(\text{output})/\partial(\text{labor})$ at fixed capital is the
\emph{marginal product of labor}. In training a neural network, every weight $w_i$ gets its own
partial $\partial(\text{loss})/\partial w_i$; the list of all of them is the gradient, and
backpropagation is an algorithm for computing millions of these partials efficiently. The
single most important fact about modern machine learning is that these partials can be computed
at all.
\end{applicationbox}

\section{Higher Partial Derivatives and Clairaut's Theorem}
\label{ch32:sec-higher}

A partial derivative $f_x$ is itself a function of $x$ and $y$, so we can differentiate it
again. This gives \emph{second-order partial derivatives}, and now there is a genuine choice:
we can differentiate first by $x$ then by $y$, or first by $y$ then by $x$. The notation, read
inner-to-outer in the operator form, is
\[
   f_{xx} = \pdv{}{x}\!\left(\pdv{f}{x}\right) = \frac{\partial^2 f}{\partial x^2},
   \qquad
   f_{yy} = \frac{\partial^2 f}{\partial y^2},
\]
\[
   f_{xy} = \pdv{}{y}\!\left(\pdv{f}{x}\right) = \frac{\partial^2 f}{\partial y\,\partial x},
   \qquad
   f_{yx} = \pdv{}{x}\!\left(\pdv{f}{y}\right) = \frac{\partial^2 f}{\partial x\,\partial y}.
\]

\begin{trapbox}
The order conventions are maddening but standard. In the \emph{subscript} form $f_{xy}$ you
read left to right: first $x$, then $y$. In the \emph{fraction} form
$\partial^2 f/\partial y\,\partial x$ you read right to left, inside out: the $\partial x$ next
to $f$ acts first. So $f_{xy}$ and $\partial^2 f/(\partial y\,\partial x)$ are the \emph{same}
object. Pick one convention and translate carefully.
\end{trapbox}

The two ``mixed'' partials $f_{xy}$ and $f_{yx}$ ask different-looking questions: how does
the eastward slope change as you move north, versus how does the northward slope change as you
move east? A priori, these rates of change need not agree. Yet under standard continuity
conditions, they are always equal.

\begin{recoverybox}
\textbf{Pause and think.}
Before calculating mixed derivatives, ask yourself what $f_{xy}$ and $f_{yx}$ physically measure.
If you stand on a hillside facing north and take a step east, your northward slope may change.
If instead you face east and step north, does your eastward slope change by the exact same amount?
It feels as though the two numbers could easily disagree on an uneven hill. Why should two
perpendicular directions always report identical twists?
\end{recoverybox}

\begin{example}[The mixed partials match]
\label{ch32:ex-mixed}
Take $f(x,y)=x^3 y^2$. Then
\[
   f_x = 3x^2 y^2, \qquad f_{xy} = \pdv{}{y}\,(3x^2 y^2) = 6x^2 y,
\]
\[
   f_y = 2x^3 y, \qquad f_{yx} = \pdv{}{x}\,(2x^3 y) = 6x^2 y.
\]
Both routes land on $6x^2 y$. Differentiating in the order $x$-then-$y$ gave the same thing as
$y$-then-$x$. Try it on $\sin(xy)$ or $e^{x}\cos y$ and the mixed partials match there too, for a
general reason.
\end{example}

\begin{theorem}[Clairaut / Schwarz]
\label{ch32:thm-clairaut}
If $f_{xy}$ and $f_{yx}$ are both continuous on an open region containing $(a,b)$, then
\[
   f_{xy}(a,b) = f_{yx}(a,b).
\]
The order of mixed partial differentiation does not matter.
\end{theorem}

\begin{intuitionbox}
Why \emph{should} order not matter? Picture a tiny rectangle with corners $(a,b)$,
$(a+h,b)$, $(a,b+k)$, $(a+h,b+k)$. Form the ``second difference'' of $f$ around the rectangle:
take the change across the top minus the change across the bottom,
\[
   \Delta = \big[f(a+h,b+k)-f(a,b+k)\big] - \big[f(a+h,b)-f(a,b)\big].
\]
Group it the other way -- change up the right side minus change up the left side -- and you get
\emph{exactly the same four terms}, just reassociated:
\[
   \Delta = \big[f(a+h,b+k)-f(a+h,b)\big] - \big[f(a,b+k)-f(a,b)\big].
\]
The first grouping, divided by $hk$ and shrunk, becomes $f_{yx}$; the second becomes $f_{xy}$.
But they were always the same number $\Delta/(hk)$. The symmetry of the mixed partials is just
the symmetry of a rectangle -- top-then-side equals side-then-top -- pushed to the limit.
\end{intuitionbox}

\begin{figure}[H]
\centering
\begin{tikzpicture}
% Left panel: Path A (horizontal differences, then vertical)
\begin{scope}[shift={(0,0)}]
  \draw[rounded corners=2pt, fill=sourceblue!6, draw=sourceblue!60, thick] (0,0) rectangle (5.0, 3.2);
  \node[font=\footnotesize\scshape, sourceblue, above] at (2.5, 3.25) {Horizontal Step First};
  
  \fill[black] (0.9,0.6) circle (2pt) node[below, font=\scriptsize] {$(a,b)$};
  \fill[black] (4.1,0.6) circle (2pt) node[below, font=\scriptsize] {$(a+h,b)$};
  \fill[black] (0.9,2.6) circle (2pt) node[above, font=\scriptsize] {$(a,b+k)$};
  \fill[black] (4.1,2.6) circle (2pt) node[above, font=\scriptsize] {$(a+h,b+k)$};
  
  \draw[->, thick, sourceblue] (1.1,0.6) -- (3.9,0.6) node[midway, below, font=\tiny] {$\Delta_x f(b)$};
  \draw[->, thick, sourceblue] (1.1,2.6) -- (3.9,2.6) node[midway, above, font=\tiny] {$\Delta_x f(b+k)$};
  \draw[->, thick, dashed, trapred] (2.5,0.9) -- (2.5,2.3) node[midway, right, font=\tiny] {Diff of $x$-diffs};
  
  \node[font=\scriptsize, align=center] at (2.5,-0.6) {
    $\Delta = \Delta_x f(b+k) - \Delta_x f(b)$\\
    $\xrightarrow{\text{MVT twice}} h\,k\,f_{yx}(d_1,d_2)$
  };
\end{scope}

% Right panel: Path B (vertical differences, then horizontal)
\begin{scope}[shift={(6.6,0)}]
  \draw[rounded corners=2pt, fill=targetgreen!6, draw=targetgreen!60, thick] (0,0) rectangle (5.3, 3.2);
  \node[font=\footnotesize\scshape, targetgreen, above] at (2.65, 3.25) {Vertical Step First};
  
  \fill[black] (1.0,0.6) circle (2pt) node[below, font=\scriptsize] {$(a,b)$};
  \fill[black] (3.8,0.6) circle (2pt) node[below, font=\scriptsize] {$(a+h,b)$};
  \fill[black] (1.0,2.6) circle (2pt) node[above, font=\scriptsize] {$(a,b+k)$};
  \fill[black] (3.8,2.6) circle (2pt) node[above, font=\scriptsize] {$(a+h,b+k)$};
  
  \draw[->, thick, targetgreen] (1.0,0.8) -- (1.0,2.4) node[midway, left, font=\tiny] {$\Delta_y f(a)$};
  \draw[->, thick, targetgreen] (3.8,0.8) -- (3.8,2.4);
  \node[targetgreen!80!black, font=\tiny, anchor=west] at (3.9,1.6) {$\Delta_y f(a+h)$};
  \draw[->, thick, dashed, trapred] (1.25,1.6) -- (3.55,1.6) node[midway, above, font=\tiny] {Diff of $y$-diffs};
  
  \node[font=\scriptsize, align=center] at (2.65,-0.6) {
    $\Delta = \Delta_y f(a+h) - \Delta_y f(a)$\\
    $\xrightarrow{\text{MVT twice}} h\,k\,f_{xy}(c_1,c_2)$
  };
\end{scope}

% Bottom connecting corridor
\begin{scope}[shift={(0,-2.3)}]
  \draw[rounded corners=2pt, fill=gray!4, draw=gray!40] (-0.1,0) rectangle (12.0, 0.9);
  \node[font=\footnotesize] at (5.95, 0.45) {
    $\Delta = f(a+h,b+k) - f(a+h,b) - f(a,b+k) + f(a,b)$ in both groupings
    $\;\Longrightarrow\; \lim_{h,k\to 0} f_{xy} = f_{yx}$.
  };
\end{scope}
\end{tikzpicture}
\caption{Why mixed partials commute: the second difference on a coordinate rectangle.
Grouping the four corner values horizontally first computes how an $x$-step changes across a $y$-interval ($f_{yx}$);
grouping vertically first computes how a $y$-step changes across an $x$-interval ($f_{xy}$).
Because both algebraic sums are identical, the Mean Value Theorem applied twice forces $f_{xy}=f_{yx}$ as the rectangle shrinks.}
\label{ch32:fig-clairaut-grid}
\end{figure}

\begin{whybox}
Here is the idea of the proof made rigorous. Define the second difference $\Delta(h,k)$ above.
Apply the one-variable Mean Value Theorem twice. Group $\Delta$ as a difference in $x$ of the
$y$-difference: the function
$\phi(x)=f(x,b+k)-f(x,b)$ satisfies $\Delta = \phi(a+h)-\phi(a) = h\,\phi'(c_1)$ for some $c_1$
between $a$ and $a+h$, and $\phi'(x)=f_x(x,b+k)-f_x(x,b)$. Apply the MVT again in the
$y$-direction to this difference: it equals $k\,f_{xy}(c_1,c_2)$ for some $c_2$ between $b$ and
$b+k$. So $\Delta = h\,k\,f_{xy}(c_1,c_2)$. Run the \emph{same} argument grouping in the other
order and you get $\Delta = h\,k\,f_{yx}(d_1,d_2)$. Therefore
$f_{xy}(c_1,c_2) = f_{yx}(d_1,d_2)$. Now let $h,k\to 0$: all four interior points are squeezed
toward $(a,b)$, and by \emph{continuity} of the mixed partials both sides converge to their
values at $(a,b)$, giving $f_{xy}(a,b)=f_{yx}(a,b)$. The continuity hypothesis is exactly what
lets the squeeze land on a single value -- without it the theorem can fail.
\end{whybox}

\begin{historybox}
Alexis Clairaut stated the symmetry of mixed derivatives around 1740 while investigating the
equilibrium shape of the spinning Earth -- a problem in celestial mechanics where gravitational
and centrifugal forces interact across several dimensions. Leonhard Euler and Jean le Rond
d'Alembert used mixed partials freely throughout the 1740s and 1750s to model vibrating musical
strings and fluid motion. At the time, mathematicians possessed only single-variable
differential calculus (Leibniz's infinitesimals) and mechanical intuition. They lacked a formal
$\eps$--$\delta$ definition of a limit, had no precise notion of two-dimensional open regions,
and blithely assumed that if partial derivatives existed along coordinate lines, they behaved
like one-variable derivatives in every respect.

That confidence led to decades of controversy. Euler and d'Alembert clashed sharply over what
curves could physically describe a plucked string: d'Alembert insisted solutions required smooth
formulas, while Euler argued that a string plucked into a triangular kink with a sharp corner
was physically real even if its derivatives were undefined at the peak. Simultaneously,
mathematicians treated the equation $\partial^2 f/\partial x\partial y = \partial^2 f/\partial y\partial x$
as an inviolable algebraic identity, believing that the order of limits could always be
interchanged. It was not until the 1870s that Hermann Schwarz, followed by Karl Weierstrass and
Giuseppe Peano, proved that symmetry can catastrophically fail if the second partials are not
continuous. Counterexamples like $f(x,y)=xy(x^2-y^2)/(x^2+y^2)$ (with $f(0,0)=0$) demonstrated
that $f_{xy}(0,0)=-1$ while $f_{yx}(0,0)=+1$, demonstrating that mixed derivative equality is a
geometric theorem about continuity rather than an algebraic birthright.

Resolving these foundations unlocked classical field theory, thermodynamics, and fluid dynamics.
In thermodynamics, physical properties like internal energy and entropy are \emph{state functions}
-- quantities that depend solely on the current state of a gas, not on how it got there. That
path-independence holds mathematically if and only if the mixed partials of the thermodynamic
potential agree. Today, multivariable calculus has scaled from two or three spatial dimensions to
millions or billions of parameters in modern deep learning. In these high-dimensional loss
functions, non-smooth activations -- such as the rectified linear unit, or ReLU, which has a sharp
hinge like an absolute-value function -- mean that second derivatives often fail to exist or
commute. Understanding optimization dynamics around saddles in non-convex multi-million-dimensional
spaces remains a vibrant, open frontier of modern mathematical research.
\end{historybox}

\begin{connectionbox}
Clairaut's theorem is the seed of a deep idea. It says the mixed partials of a smooth function
agree -- equivalently, that a gradient vector field $(\,f_x, f_y\,)$ (an assignment of a little
arrow of slopes to every point in the plane) has a special symmetry:
$\partial_y f_x = \partial_x f_y$. In vector calculus this becomes the statement
``the curl of a gradient is zero,'' where \emph{curl} measures the tendency of a fluid flow to
spin a tiny paddle wheel placed at that point. A vector field whose curl vanishes is
\emph{conservative} -- meaning it represents a force field (like gravity) where moving a particle
around any closed loop costs zero net energy, precisely because the force derives from an
underlying potential energy function. That equality $f_{xy}=f_{yx}$ returns in Chapter~35 as the
test for whether a planar force field comes from a potential function.
\end{connectionbox}

\section{Tangent Planes}
\label{ch32:sec-tangent}

In one variable, the tangent line was the best straight-line approximation to a curve: zoom in
on a smooth graph and it looks linear. The same zooming works on a surface, but a flat
approximation in the plane is not a line -- it is a \emph{plane}. The tangent plane is the best
flat approximation to the surface at a point, and the partial derivatives are precisely the two
slopes that pin it down.

\begin{intuitionbox}
Stand on the surface at the point above $(a,b)$. You found two slopes there: $f_x(a,b)$, the
tilt as you face east, and $f_y(a,b)$, the tilt as you face north. A flat plane through that
point is completely determined by how steeply it rises eastward and how steeply it rises
northward. Match those two slopes to the surface's two slopes and you have the unique plane that
approximates the surface -- the tangent plane. East-slope and north-slope are exactly enough data to
fix a plane, which is why \emph{two} partials suffice.
\end{intuitionbox}

A plane through the point $(a,b,f(a,b))$ has the form
\[
   z = f(a,b) + A\,(x-a) + B\,(y-b),
\]
where $A$ is its slope in the $x$-direction and $B$ its slope in the $y$-direction. To be
tangent, those slopes must equal the surface's slopes: $A=f_x(a,b)$ and $B=f_y(a,b)$. That
gives the formula.

\begin{focusbox}
\textbf{Tangent plane.} If $f$ has continuous partials near $(a,b)$, the tangent plane to
$z=f(x,y)$ at the point $\big(a,b,f(a,b)\big)$ is
\[
   \boxed{\,z = f(a,b) + f_x(a,b)\,(x-a) + f_y(a,b)\,(y-b)\,}.
\]
The right-hand side is the \emph{linearization} $L(x,y)$ of $f$ at $(a,b)$: the best
linear approximation, just as the tangent line was in one variable.
\end{focusbox}

\begin{happeningbox}
Compare this with the one-variable tangent line $y = f(a) + f'(a)(x-a)$ -- which we saw in
Chapter~31 as the first-order Taylor approximation $P_1(x)$. The structure is identical:
value at the base point, plus slope times displacement. We simply have \emph{two} displacement
terms now, one per input direction, each weighted by its own partial derivative. The tangent
plane is two tangent lines -- the slice-curves of Figure~\ref{ch32:fig-slices} -- spanned into a
flat sheet. ``Differentiable at $(a,b)$'' means that this plane really is an excellent
approximation: the surface peels away from the plane faster than you walk, so the error shrinks
faster than the distance $\sqrt{(x-a)^2+(y-b)^2}$. (Beware: in several variables the mere
\emph{existence} of $f_x$ and $f_y$ is not enough to guarantee this -- continuity of the partials
is the clean condition that does.)
\end{happeningbox}

\begin{figure}[H]
\centering
\begin{tikzpicture}
\begin{axis}[
title style={font=\small},
  width=10.5cm, height=7.4cm,
  view={38}{20},
  xlabel={$x$}, ylabel={$y$}, zlabel={$z$},
  domain=-0.4:2.4, y domain=-0.9:1.9,
  samples=24, samples y=24,
  colormap/blackwhite,
  title={The tangent plane to $z=x^2+y^2$ at $(1,0.5)$},
  zmin=-0.5, zmax=7,
]
\addplot3[surf, opacity=0.30, shader=flat, draw=gray!45] {x^2 + y^2};
% tangent plane: z = 1.25 + 2(x-1) + 1(y-0.5)
\addplot3[
  draw=tangentred, fill=tangentred!18, opacity=0.55,
  samples=2, samples y=2, domain=0.0:2.0, y domain=-0.5:1.5,
] {1.25 + 2*(x-1) + 1*(y-0.5)};
\addplot3[mark=*, mark size=1.7pt, black] coordinates {(1,0.5,1.25)};
\end{axis}
\end{tikzpicture}
\caption{The bowl and its tangent plane at $(1,0.5,1.25)$. The plane has $x$-slope
$f_x=2$ and $y$-slope $f_y=1$, matching the two slice-slopes from
Figure~\ref{ch32:fig-slices}. It touches the surface at one point and is the surface's best
linear approximation nearby.}
\label{ch32:fig-tangentplane}
\end{figure}

\begin{example}[Building a tangent plane]
\label{ch32:ex-tangent}
Find the tangent plane to $f(x,y)=x^2+y^2$ at $(1,0.5)$.

We already have $f(1,0.5)=1+0.25=1.25$, and from the lab box $f_x(1,0.5)=2$,
$f_y(1,0.5)=1$. Plug in:
\[
   z = 1.25 + 2(x-1) + 1\,(y-0.5) = 2x + y - 1.25 .
\]
That is the plane in Figure~\ref{ch32:fig-tangentplane}.
\end{example}

\begin{example}[Linear approximation in action]
\label{ch32:ex-linapprox}
Use the tangent plane to estimate $f(1.1, 0.45)$ for $f=x^2+y^2$ without squaring anything.
The linearization is $L(x,y)=2x+y-1.25$, so
\[
   L(1.1,0.45) = 2(1.1) + 0.45 - 1.25 = 2.2 + 0.45 - 1.25 = 1.40 .
\]
The true value is $1.1^2 + 0.45^2 = 1.21 + 0.2025 = 1.4125$. The linear estimate $1.40$ is off
by only $0.0125$ -- excellent for a step of about a tenth, and a perfect echo of how tangent
lines estimated $\sqrt{4.1}$ back in the one-variable chapter. The whole machinery of
differentials and linear approximation survives the jump to several variables; it just acquires
a second slope term.
\end{example}

\begin{aestheticbox}
The linear approximation generalizes directly across dimensions. In Chapter~31, Taylor's theorem began with the degree-one
linear polynomial: $P_1(x) = f(a)+f'(a)(x-a)$. In two variables, that linear approximation is
the tangent plane: $L(x,y) = f(a,b)+f_x(a,b)\,(x-a)+f_y(a,b)\,(y-b)$. In three variables, it is
a tangent hyperplane: $f(a,b,c)+f_x(x-a)+f_y(y-b)+f_z(z-c)$. The pattern continues consistently,
adding one linear slope term per dimension.

Knowing rates of change due east ($f_x$) and due north ($f_y$) raises a natural question:
what about arbitrary directions? In Chapter~33, these separate partial derivatives assemble into
the \emph{gradient vector} $\grad f = (f_x, f_y, f_z)$. With the gradient, directional
derivatives become simple dot products, the chain rule extends to paths, and linearization
takes the vector form $f(\vec a) + \grad f(\vec a)\cdot(\vec x - \vec a)$.
\end{aestheticbox}

\section{Exercises}
\label{ch32:sec-exercises}

\begin{exercise}[Warm-up mental checks]
\label{ch32:exc-mental-checks}
Test your immediate geometric and algebraic intuition with these quick mental checks (no scratch paper required):
\begin{enumerate}[label=(\alph*)]
  \item At the point $(x,y)=(1,2)$, what is $\pdv{}{x}\,(x^3 y^4)$? (Hold $y=2$ fixed as a constant number; differentiate $16x^3$ with respect to $x$).
  \item A hiker on a pass discovers that moving due east has zero slope ($\pdv{z}{x}=0$) and moving due north has zero slope ($\pdv{z}{y}=0$). Is the hiker guaranteed to be standing at a mountain peak or valley? (Think of the saddle $z=x^2-y^2$).
  \item Suppose a classmate verifies that $\lim_{x\to 0} f(x, mx) = 0$ along every straight line of slope $m$. Can they conclude that $\lim_{(x,y)\to(0,0)} f(x,y) = 0$? Which counterexample curve from this chapter refutes it?
  \item If $f(x,y) = \sin(x) e^y$, why must $\pdv{^2 f}{x\,\partial y} = \pdv{^2 f}{y\,\partial x}$ without writing down any derivatives? Which theorem guarantees this?
\end{enumerate}
\end{exercise}

\begin{exercise}[Domains and contours]
\label{ch32:exc-domain}
For each function, describe the domain as a region in the plane and sketch three level curves.
\begin{enumerate}[label=(\alph*)]
  \item $f(x,y)=\sqrt{9-x^2-y^2}$.
  \item $f(x,y)=\ln(x-y)$.
  \item $f(x,y)=\dfrac{1}{x^2+y^2}$ (what happens to the contours near the origin?).
\end{enumerate}
\end{exercise}

\begin{exercise}[Reading a saddle]
\label{ch32:exc-saddle}
For $f(x,y)=y^2-x^2$, identify the level curve at $c=0$, at $c=1$, and at $c=-1$. Which way does
the surface rise from the origin, and which way does it fall? Sketch the contour map and mark the
saddle point.
\end{exercise}

\begin{exercise}[The two-path test]
\label{ch32:exc-twopath}
Show that each limit fails to exist by finding two paths giving different values.
\begin{enumerate}[label=(\alph*)]
  \item $\displaystyle\lim_{(x,y)\to(0,0)} \frac{x^2-y^2}{x^2+y^2}$.
  \item $\displaystyle\lim_{(x,y)\to(0,0)} \frac{x^2 y}{x^4+y^2}$ (try lines, then the parabola
  $y=x^2$).
\end{enumerate}
\end{exercise}

\begin{exercise}[A limit that exists]
\label{ch32:exc-exists}
Using polar coordinates or a squeeze, prove that
\[
   \lim_{(x,y)\to(0,0)} \frac{x^3 - y^3}{x^2+y^2} = 0.
\]
Then explain in one sentence why the angle $\theta$ disappearing from your bound is exactly what ensures the limit exists.
\end{exercise}

\begin{exercise}[Computing partials]
\label{ch32:exc-partials}
Find $f_x$ and $f_y$.
\begin{enumerate}[label=(\alph*)]
  \item $f(x,y)=x^4 y - 3xy^3 + 7$.
  \item $f(x,y)=e^{xy}\sin x$.
  \item $f(x,y)=\dfrac{x}{x^2+y^2}$.
  \item $f(x,y)=\ln\!\big(x^2 + y^2\big)$.
\end{enumerate}
\end{exercise}

\begin{exercise}[Three variables]
\label{ch32:exc-three}
For $g(x,y,z) = x\,y^2\,z^3 + \dfrac{x}{y} - \cos(yz)$, compute all three first partials
$g_x$, $g_y$, $g_z$.
\end{exercise}

\begin{exercise}[Verifying Clairaut]
\label{ch32:exc-clairaut}
For $f(x,y)=x^2 e^{y} + \sin(xy)$, compute $f_{xy}$ and $f_{yx}$ separately and confirm they
agree. State explicitly where in your work you used the fact that the function is built from
continuous pieces.
\end{exercise}

\begin{exercise}[Tangent plane and estimate]
\label{ch32:exc-tangent}
Let $f(x,y)=\sqrt{x^2+y^2}$ (the cone).
\begin{enumerate}[label=(\alph*)]
  \item Find the tangent plane at $(3,4)$, where $f=5$.
  \item Use it to estimate $f(3.1, 3.9)$, and compare with the exact value.
  \item Explain why no tangent plane exists at the origin. (Hint: walk toward $(0,0)$ from two
  directions and look at the slopes -- this is the cone's sharp tip.)
\end{enumerate}
\end{exercise}

\begin{exercise}[Conceptual: lines are not enough]
\label{ch32:exc-concept}
A student claims: ``I checked that
\[
   \lim_{(x,y)\to(0,0)} \frac{x^2 y}{x^4+y^2} = 0
\]
along every straight line $y=mx$, so the multivariable limit must be $0$.''
In a short paragraph, explain what is wrong with this conclusion, and state the curved path that refutes it.
\end{exercise}

\begin{exercise}[Proof: Clairaut by the second difference]
\label{ch32:exc-proof}
Let $\Delta(h,k) = f(a+h,b+k) - f(a+h,b) - f(a,b+k) + f(a,b)$. Assuming $f$, $f_x$, $f_y$,
$f_{xy}$, and $f_{yx}$ are all continuous near $(a,b)$, apply the Mean Value Theorem twice in
each of the two groupings of $\Delta$ to show $f_{xy}(a,b)=f_{yx}(a,b)$. Identify the precise
step where continuity of the mixed partials is needed.
\end{exercise}

\begin{exercise}[A function continuous along every line but not at the origin]
\label{ch32:exc-pathology}
Let $g(x,y)=\dfrac{x^2 y}{x^4+y^2}$ for $(x,y)\neq(0,0)$ and $g(0,0)=0$.
\begin{enumerate}[label=(\alph*)]
  \item Show $g$ restricted to any line through the origin is continuous at the origin.
  \item Show $g$ is nonetheless \emph{not} continuous at the origin.
  \item In one sentence, say what this example teaches about
  Definition~\ref{ch32:def-continuity}.
\end{enumerate}
\end{exercise}
